\chapter{Validación y evaluación}
\label{chap:validacion}

\lettrine{L}{a} validación se orienta a comprobar la viabilidad del sistema como arquitectura operativa y a medir el grado de madurez de integración entre sus capas.

\section{Metodología de validación}

La evaluación se diseña con enfoque de ingeniería aplicada, separando tres niveles:

\begin{itemize}
  \item \textbf{validación funcional}: cada bloque cumple su responsabilidad;
  \item \textbf{validación de integración}: interoperabilidad entre Android, ESP32 y repositorio de sesiones;
  \item \textbf{validación operativa}: estabilidad del flujo de trabajo en escenarios de uso repetido.
\end{itemize}

El objetivo no es obtener un único resultado numérico, sino demostrar que el sistema funciona de forma trazable y mantenible.

\section{Escenarios de prueba}

Los escenarios principales considerados son:

\begin{itemize}
  \item captura y gestión de sesión en modo móvil;
  \item operación remota local sobre ESP32 (estado, inicio/parada y descarga);
  \item consistencia del contrato de sesión entre orígenes móvil y embebido;
  \item robustez de operación tras múltiples iteraciones de captura y consulta;
  \item ensayo preliminar de capacidad de comunicación en escenario multi-nodo, para estimar el rango operativo del sistema ante concurrencia de varios emisores.
\end{itemize}

\figplaceholder{%
Diagrama recomendado:\\
escenarios de validación por subsistemas y flujo de integración
}{Escenarios de prueba del sistema.}
\label{fig:placeholder-protocolo-experimental}

\section{Configuración de pruebas}

Para cada prueba se documentan, como mínimo:

\begin{itemize}
  \item nodo origen de los datos (móvil o embebido);
  \item parámetros de captura usados en la sesión;
  \item versión de app/firmware empleada;
  \item ruta de comunicación utilizada (local o sincronizada);
  \item resultado de operación (éxito, error o comportamiento degradado);
  \item número de nodos concurrentes, tasa de envío y tamaño de carga útil en las pruebas de capacidad.
\end{itemize}

Esta documentación permite reproducibilidad y facilita análisis de incidencias.

\section{Métricas de evaluación}

Las métricas se centran en comportamiento del sistema:

\begin{itemize}
  \item tasa de sesiones completadas sin error;
  \item latencia percibida en operaciones de control (start/stop/descarga);
  \item consistencia de metadatos y formato de sesión;
  \item estabilidad de conexión en modo embebido;
  \item esfuerzo operativo (pasos necesarios por flujo);
  \item degradación de latencia o pérdida de operaciones al aumentar el número de nodos concurrentes.
\end{itemize}

\section{Plantilla de comparativa móvil vs embebido}

\begin{table}[htbp]
  \centering
  \small
  \setlength{\tabcolsep}{6pt}
  \renewcommand{\arraystretch}{1.2}
  \rowcolors{2}{white}{udcgray!25}
  \begin{tabularx}{\textwidth}{>{\raggedright\arraybackslash}p{0.2\textwidth} >{\raggedright\arraybackslash}X >{\raggedright\arraybackslash}X >{\raggedright\arraybackslash}p{0.24\textwidth}}
  \rowcolor{udcpink!25}
  \textbf{Aspecto} & \textbf{Nodo móvil (Android)} & \textbf{Nodo embebido (ESP32+IMU)} & \textbf{Criterio de evaluación} \\\hline
  Flujo de captura & Local y directo & Remoto con control desde app & Completitud de flujo \\
  Persistencia & Sesión local en app & Descarga y guardado en app & Trazabilidad de sesión \\
  Operación & Interacción simple & Mayor dependencia de integración & Robustez operativa \\
  Estado/diagnóstico & Inmediato en interfaz & Requiere chequeo de conectividad & Tiempo de resolución \\
  Evolución cloud & Integración prevista con backend & Integración prevista vía contratos comunes & Escalabilidad de arquitectura \\
  \end{tabularx}
  \caption{Plantilla de comparación entre modo móvil y modo embebido.}
  \label{tab:comparativa-plataformas}
\end{table}

\section{Criterios de aceptación}

Se considerará validada la arquitectura en el alcance del TFG cuando:

\begin{itemize}
  \item el flujo captura-guardado-consulta funcione de forma repetible;
  \item la integración Android-ESP32 permita control y descarga de sesiones;
  \item el formato de sesión sea consistente entre orígenes;
  \item se documenten límites operativos y decisiones de evolución pendientes.
\end{itemize}

\section{Resultados}

Este apartado se reserva para presentar resultados reales de campaña, distinguiendo claramente entre validación funcional, integración y operación.

\subsection{Resultados del modo móvil}

Incluir sesiones representativas de captura local, tiempos de operación y comportamiento de gestión de sesiones.

\figplaceholder{%
Resultado pendiente:\\
captura de flujo de sesión en modo móvil
}{Resultados representativos del modo móvil.}
\label{fig:placeholder-resultado-movil}

\subsection{Resultados del modo embebido}

Documentar resultados de conexión, control de captura y descarga de sesiones desde ESP32, incluyendo incidencias observadas.

\figplaceholder{%
Resultado pendiente:\\
captura de estado y sesión en modo embebido
}{Resultados representativos del modo embebido.}
\label{fig:placeholder-resultado-embebido}

\subsection{Evaluación del subsistema de comunicación}

Este bloque debe recoger la comparación entre los mecanismos de comunicación considerados para el enlace Android-ESP32. La evaluación debe centrarse en latencia, estabilidad de operación, simplicidad de integración y comportamiento cuando aumenta la carga de mensajes o el número de nodos concurrentes.

Como guía de redacción, conviene documentar:

\begin{itemize}
  \item tecnologías comparadas;
  \item tamaño de mensaje y frecuencia de envío usados en la prueba;
  \item tiempos de respuesta observados en operaciones de control;
  \item pérdidas, errores o degradación bajo carga;
  \item criterio seguido para seleccionar la opción más adecuada.
\end{itemize}

\figplaceholder{%
Resultado pendiente:\\
comparativa de latencia, estabilidad y carga entre mecanismos de comunicación
}{Evaluación comparativa del subsistema de comunicación Android-ESP32.}
\label{fig:placeholder-resultado-comunicacion}

\subsection{Resultados de integración por contratos}

Este bloque debe mostrar si los contratos de sesión y control se cumplen de forma homogénea y dónde aparecen incompatibilidades.

\figplaceholder{%
Figura recomendada:\\
resumen de consistencia de contratos entre nodos
}{Validación de contratos de integración entre subsistemas.}
\label{fig:placeholder-resultado-espectral}

\subsection{Comparación de operación entre modos}

La comparación debe centrarse en esfuerzo operativo, estabilidad de flujo y mantenibilidad, no solo en valores puntuales.

\figplaceholder{%
Figura recomendada:\\
comparativa operativa móvil vs embebido
}{Comparación operativa entre modo móvil y modo embebido.}
\label{fig:placeholder-comparacion-final}

\section{Discusión}

La discusión final debe abordar:

\begin{itemize}
  \item fortalezas reales del diseño modular;
  \item puntos de fricción en integración embebida;
  \item impacto de no tener todavía backend cloud en producción;
  \item criterio técnico para decidir la evolución de mecanismos de comunicación y backend.
\end{itemize}

\subsection{Limitaciones del trabajo}

Las limitaciones deben expresarse en términos de alcance de implementación, cobertura de pruebas y grado de madurez de integración.

\subsection{Lecciones aprendidas}

Como cierre de validación, este apartado debe recoger decisiones que hayan reducido riesgo técnico y pautas que mejoren la siguiente iteración del sistema.
